% ============================================================================
%  RESULTS -- MAIN TEXT, SECTION 4: ASYMPTOTIC PROPERTIES
%  project `panel-npiv-discrete`
%
%  WHAT THIS FILE HOLDS.  The 4 statements that belong under Section 4:
%  the four asymptotic statements: the second-order remainder on the ball, the
%  rate for the varphi-step, the root-n theorem, and the rate floor.
%  Nothing else.  In order: prp:sharpquad, prp:phirate, thm:rootn, cor:ratefloor.
%
%  WHY THREE FILES AND NOT ONE.  results_main.tex printed all 23 core statements
%  in one block, which the article had to \input somewhere; it went under
%  Section 4, so Section 2's prose forward-referenced statements forty pages
%  later.  The three files are \input under the three sections that discuss
%  them.  The split is a PARTITION: results_s2 + results_s3 + results_s4 +
%  results_supp = the undivided results.tex, every \label declared exactly once
%  across the four, no statement rewritten, retitled or renumbered.  Relative
%  order inside each file is the order of the parent file.
%
%%  DO NOT \input THIS FILE ALONGSIDE results.tex OR assumptions.tex.
%%  results.tex and assumptions.tex SURVIVE UNDIVIDED and remain the source of
%%  record: scripts/build-theory/build.sh compiles them, scripts/render-sync/
%%  check_rendering_sync.py reads them, and ~100 files under scripts/symbolic/
%%  specs/ cite them BY LINE NUMBER.  Neither was edited.  Every \label in this
%%  file is therefore declared TWICE on disk -- here and in its parent -- and the
%%  two must never reach the same document.
%
%  SUPPLEMENT POINTERS.  A statement kept here may cite a statement that lives
%  in the supplement.  Those calls are written \suppref{Kind}{name} ->
%  "Kind S:name" and \suppeq{name} -> "(S:name)", both defined in
%  paper_header.tex.  They are not \cref calls and cannot render "??".
% ============================================================================


% THE THREE DEFINITIONS BELOW WERE ADDED 2026-08-05 and sit AHEAD of the
% proposition rather than inside it.  B_sym, q and sigma_q appear in the
% proposition's own precis and in its numbered conclusions eq:sharpconst and
% eq:bracket, and were defined only in the supplement -- so the article printed a
% sharp-constant equality and a bracketing whose norms had no antecedent in the
% document a reader was holding.  A first attempt put them after the precis,
% which left sigma_q still used before it was defined, inside the same printed
% environment; setup belongs before the statement it is setup for.
Three objects in what follows are constructed in the supplement and are
recalled here so that the statement can be read where it is printed.
$\mathcal B_{\mathrm{sym}}:=\tfrac12(\mathcal B+\mathcal B^{\ast})$ is the
symmetric part of $\mathcal B$; $q:=E[\partial_u\ell_\star\,\widetilde W\mid
X_1,X_2]\in\IR^{p}$, the $\IR^{p\times(\mathsf L-1)}$-valued $\partial_u\ell_\star$
contracted against the contrast vector $\widetilde W$ of \eqref{eq:contrast} and then
conditioned; and, for an integrable $g$, $\sigma_g(x):=E[g\mid X_2=x]f_{X_2}(x)
+E[g\mid X_1=x]f_{X_1}(x)$, taken coordinatewise for a vector-valued $g$, so that
$\sigma_q$ is that object at $g=q$ and is $\IR^{p}$-valued. The
multiplier is written out as $\vartheta'\sigma_q/\pi$ throughout and is given no
abbreviation, because $\mathcal B_{\mathrm{mult}}$ is multiplication \emph{by} it
and an abbreviation would collapse the distinction.
At $\mathsf L=2$ the contraction is a product of scalars in the second factor,
$\partial_u\ell_\star\widetilde W=\widetilde W\partial_u\ell_\star$, and $q$, $\sigma_q$ and
$\vartheta'\sigma_q/\pi$ are the objects earlier rounds printed.

\begin{proposition}[The second-order remainder on the ball: the sharp constant, and the exact $\delta^{2}$ order under O-16(a)]\label{prp:sharpquad}
The supremum of the second-order remainder over a ball of radius $\delta$ in the structural direction is
exactly $\tfrac12\delta^{2}$ times the norm of the compressed symmetrised form, and it is bounded below by
$\tfrac12\delta^{2}\|\vartheta'\sigma_q/\pi\|_\infty$; so, provided the form does not degenerate, no bound
of the form $o(\delta^{2})$, and in particular no bound in the weak norm
$\|\mathcal A_{\mathrm{cond}}\cdot\|_{L^{2}(\Pi)}$, holds uniformly on that ball.
Suppose the hypotheses of \suppref{Lemma}{solset} hold, and for clause \textup{(iv)} in addition
\textbf{C2$'$(i)}. Fix $\vartheta\in\IR^{p}$ with $\|\vartheta\|=1$,
and let the \emph{second-order remainder in the structural direction} at $h\in\mathcal V$ be
\begin{equation}\label{eq:quadrem}
\tfrac12\,\vartheta'\partial^{2}_{\tau_1\tau_1}\Psi(0)
\;=\;\tfrac12\,E\big[\vartheta'\,\partial_u\ell_\star\,\widetilde W\,(\Delta h)^{2}\big]
\;=\;\tfrac12\,\big\langle\mathcal Bh,h\big\rangle_{L^{2}(\pi)} ,
\end{equation}
which is \suppref{Lemma}{secondord}\textup{(i)}, read at the admissible direction pair
$(D,\mathsf r)=(0,0)$, contracted in $\vartheta$ and transferred to $L^{2}(\pi)$ by
\suppref{Lemma}{solset}\textup{(iii)}.
% WHY lem:solset IS CITED BY CLAUSE, what reading lem:secondord at (0,0) buys and
% what it does not inherit, and why the remainder carries no symbol of its own:
% rem:sharpquadread.
\emph{Two non-degeneracy hypotheses are carried by name below, and this display is where they are
defined}: at the fixed $\vartheta$ above,
% DEFINITION RESTORED 2026-08-27.  O-16(a) and O-16(b) were carried as NAMED
% HYPOTHESES in the title, in (ii), in (iii), in (iv), in rem:sharpquadread(1),
% in rem:sharpscope(3)-(4), in rem:rootnscope and at four sites in prp:sharpquad's
% proof -- and were DEFINED NOWHERE in the compiled corpus.  The undivided mirror
% defined (a) as an EM-DASH APPOSITION inside clause (iii) ("Under O-16(a) --- there
% exists h in V with vartheta' E[q (Delta h)^2] != 0 --- the constant ..."), and the
% split replaced the dashes with COMMAS, turning a definition into a comma-chain of
% consequences.  Restored here as a display, once, before the clauses that read it,
% rather than re-inserting the dashes: a hypothesis carried by name across three
% results and two remarks is not an apposition inside one clause of one of them.
%   THE DISPLAY IS DELIBERATELY UNNUMBERED AND CARRIES NO \label.  A numbered
% display here would advance the ARTICLE's equation counter and renumber every
% display after it in Secs. 4-5; and it would be referenced by nobody, since all
% nine consumer sites already cite these hypotheses BY NAME ("under O-16(b)"),
% which is the reference mechanism the tags exist to provide.  C110a's rule --
% count the counter advance before taking a label -- applied in the direction of
% declining one.
\[
\text{\textup{O-16(a)}}:\quad \exists\,h\in\mathcal V\ \text{ with }\ \vartheta'E\big[q\,(\Delta h)^{2}\big]\neq0 ;
\qquad\qquad
\text{\textup{O-16(b)}}:\quad \vartheta'\sigma_q\not\equiv0 .
\]
\emph{Both are hypotheses and neither is a consequence of O-11}, which is
\suppref{Remark}{sharpscope}\textup{(3)}--\textup{(4)}; \emph{that \textup{(b)} implies \textup{(a)},
and that the converse fails, are claims of clause \textup{(iii)} below and of
\suppref{Remark}{sharpscope}\textup{(3)} respectively, not background assumed here}.
Then, for every $\delta>0$:
\begin{enumerate}
\item[(i)] \emph{The sharp constant.}
\begin{equation}\label{eq:sharpconst}
\sup_{h\in\mathcal V,\ \|h\|_{L^{2}(\pi)}\le\delta}\ \tfrac12\,\big|\big\langle\mathcal Bh,h\big\rangle_{L^{2}(\pi)}\big|
\;=\;\tfrac12\,\delta^{2}\,\big\|P_{\mathcal V}\mathcal B_{\mathrm{sym}}P_{\mathcal V}|_{\mathcal V}\big\| ,
\end{equation}
a supremum that need not be attained.
\item[(ii)] \emph{Bracketing.} The constant in \eqref{eq:sharpconst} obeys
\begin{equation}\label{eq:bracket}
\big\|\vartheta'\sigma_q/\pi\big\|_\infty
\;\le\;\big\|P_{\mathcal V}\mathcal B_{\mathrm{sym}}P_{\mathcal V}|_{\mathcal V}\big\|
\;\le\;\big\|\mathcal B_{\mathrm{sym}}\big\|\;\le\;\big\|\mathcal B\big\| ,
\end{equation}
where the first inequality holds under O-16\textup{(b)} and the outer two hold always.
% The left-hand member is a LOWER bound and the right-hand members are UPPER
% bounds; neither is the sharp constant.  rem:sharpquadread(4).
\item[(iii)] Positivity, which is a hypothesis and is stated as one. Under
\textbf{O-16\textup{(a)}} the constant in \eqref{eq:sharpconst} is strictly positive, so
$\sup_{\|h\|\le\delta}\tfrac12|\langle\mathcal Bh,h\rangle|\asymp\delta^{2}$ with a constant free of
$\delta$, and no bound of the form $o(\delta^{2})$ as $\delta\downarrow0$ holds uniformly over
$\{h\in\mathcal V:\|h\|_{L^{2}(\pi)}\le\delta\}$. Under \textbf{O-16\textup{(b)}} the constant is at
least $\|\vartheta'\sigma_q/\pi\|_\infty>0$, and O-16\textup{(b)} implies O-16\textup{(a)}.
\item[(iv)] \emph{No weak-norm bound.} Under O-16\textup{(b)} there is no finite $\lambda$ with
\begin{equation}\label{eq:noweak}
\tfrac12\big|\big\langle\mathcal Bh,h\big\rangle_{L^{2}(\pi)}\big|
\;\le\;\lambda\,\big\|\mathcal A_{\mathrm{cond}}h\big\|^{2}_{L^{2}(\Pi)}
\qquad\text{for all }h\in\mathcal V ;
\end{equation}
nor with $\|\mathcal A_{\mathrm{cond}}h\|^{2}$ replaced by
$\|h\|_{L^{2}(\pi)}\,\|\mathcal A^\ast_{\mathrm{cond}}\mathcal A_{\mathrm{cond}}h\|_{L^{2}(\pi)}$.
\end{enumerate}
\end{proposition}

\noindent\emph{Notes on \cref{prp:sharpquad}.} How its hypothesis list is read is
\suppref{Remark}{sharpquadread}; what it claims and the four things it does not are
\suppref{Remark}{sharpscope}.

% The internal label was printed in this heading through dispatch 25: the article
% rendered "Proposition 25 (prp:phirate: the rate for the varphi-step)".  A label
% is a source-side handle, not part of the result's name.
\begin{proposition}[The rate for the $\varphi$-step]\label{prp:phirate}
Suppose A0, A1, A1b(iii)--(iv), A4(ii)--(iii), A5(vi$'$), B1(i)--(v), C1(ii), C2$'$, C3 and C5$'$(i) hold.
% A4, C1 AND C5' ARE CITED BY CLAUSE AND NOT AS A BLOCK, and each restriction drops
% a clause this proposition cannot have.  Why, with the parameter configuration at
% which a block citation of C5' would make the hypothesis set empty:
% rem:phirhyp(1).
% A4 NARROWED 2026-08-10 (Phase 6a).  The tag read "A4" through dispatch 25 while
% the comment two lines below announced this proposition's own by-clause policy for
% C1 and C5', so the line applied its stated rule to two tags out of three.
% Sweeping the whole proof: A4 is cited at two steps and both are A4(ii)--(iii);
% A4(i) is consumed at NO step.  The downstream disclaimers are already explicit --
% lem:unifop prints that eq:wbounds "is not on this line: that display is two-sided
% and its lower half rests on A4(i), which neither appeal reads" -- so the block
% citation contradicted a result this proposition consumes.  NOTHING INHERITS THE
% NARROWING BY SURPRISE: thm:rootn's (T-5) enumerates what it takes from here and
% never names A4, and (T-3) carries A4(i), (ii), (iii) on its own line.
Suppose H1 holds, \emph{for the Sobolev scale $H^{s}$ that \suppref{Definition}{srcball}, A5(vi$'$),
(P-c) and (P-g) are written in, and for A2's Picard form; the identity $\mu=s/(2a)$ is \textbf{not}
assumed and is \textbf{not} used, $\mu$ being primitive at (P-a.3)}. Suppose A8 holds in its
\emph{basis} clause, $b_1,\dots,b_{L_x}$ is a B-spline basis of degree exceeding $s-\tfrac12$ on a
quasi-uniform mesh over $\mathcal X$, and in its \emph{upper end} $L_x\ll n^{1/2}$, but \textbf{not} in
its lower end $L_x\gg n^{1/(2s)}$; and suppose $p_K\ge s-\tfrac12$, which with
$p_K=8$ caps $s\le\tfrac{17}2$ and therefore, against \eqref{eq:lxwindow}'s floor,
\emph{restricts $\alpha<\tfrac{68}{145}$}.
% THE alpha RESTRICTION IS PRINTED HERE AS OF 2026-08-05.  It was previously
% derivable only by combining two clauses and was worked out in rem:phirhyp(2),
% a supplement remark; the statement itself imposed nothing on alpha, so a reader
% checking the hypotheses at r = 8, d_Z = 1 (alpha = 8/17, frontier and eq:muhalf
% both clear) would find no s and no warning.  On that band this proposition and
% thm:rootn's group (T-5) are vacuous rather than false.  A hypothesis a reader
% can check is not the same thing as a remark that records the emptiness.
% ARCHAEOLOGY.  Through dispatch 20 the p_K clause was said to be "free at every s
% the budget table admits", which is FALSE: with p_K = 8 fixed by C3 and C5'(i) the
% clause is a hard cap s <= 17/2, and it collides with eq:lxwindow's floor at every
% alpha >= 68/145.  The disposition -- vacuity, not falsity -- is rem:phirhyp(2).
And suppose
\begin{enumerate}
\item[(P-a)] \emph{(preliminary, and the \textbf{spectral} source condition at the compression)} Three
clauses, all three \emph{envelopes holding with probability tending to one} rather than $O_P$
statements, \suppref{Remark}{paenvelope} says why the distinction is load-bearing.
\begin{enumerate}
\item[(P-a.1)] \emph{(preliminary, $H^{s}$, and $\sigma(\mathcal I_\ast)$-measurable)}
$\hat\varphi^{(0)}$ is a $\sigma(\mathcal I_\ast)$-measurable random element of $\mathcal V$, in the sense and
with the fold $\mathcal I_\ast$ of \cref{def:lfstep}; $\hat\varphi^{(0)}\in\mathcal S_n$; and
$P\big(\|\vdagger-\hat\varphi^{(0)}\|_{H^{s}}\le\mathsf R\big)\to1$, with $\mathsf R$ the
\emph{same} deterministic radius as in \suppref{Definition}{srcball}, A5(vi$'$), (P-c) and (P-g).
\emph{A5(vi$'$), (P-c) and (P-g) are named here as instances and this clause maintains no list}:
\suppref{Definition}{srcball} carries the enumeration --- seven sites, this clause itself the third of
them --- and this points at it. \emph{The distinction is not clerical here, because this clause is a
\textbf{hypothesis}} and a reader must act on it: a radius chosen to clear only the sites named above
need not clear \suppeq{entropy} or \suppref{Lemma}{unifop} Step~4, which are the remaining two, and the
existential is over \textbf{one} $\mathsf R$ meeting all seven at once.
% WHY THE MEASURABILITY HALF IS STATED HERE, and why the two remaining halves are
% jointly satisfiable only under eq:sieveinhs: rem:phirhyp(3).
% "THIRD BULLET" WAS RETIRED 2026-08-02 AND THE REASON IS STRUCTURAL, NOT STYLISTIC.
% def:lfstep's measurability verification was four bullets when this sentence was
% written; the triage of that day moved them to results_s3_supp.tex, so in the
% ARTICLE the phrase named a bullet the article no longer prints -- a printed
% hypothesis pointing at absent text.  "Measurability paragraph" names what the
% article DOES retain and stays true on either side of the split.
% NOT CHANGED IN results.tex: that file is the UNDIVIDED corpus source, still
% carries all eight bullets, and "third bullet" remains correct there.  The two
% files disagree because the documents disagree, which is the point of the split.
% ARCHAEOLOGY.  Through dispatch 19 the measurability half was stated only inside
% that verification -- a hypothesis a proof reads that no hypothesis list stated.
\item[(P-a.2)] \emph{(preliminary, rate)} $P\big(\|\vdagger-\hat\varphi^{(0)}\|_{L^{2}(\pi)}\le\hat\varepsilon_0\big)\to1$
for a known deterministic $\hat\varepsilon_0=O(n^{-\varsigma_0})$ whose exponent clears the repelling
fixed point of the stagewise recursion,
\[
\varsigma_0\;>\;\frac{\kappa\,\alpha}{4\,(2\kappa-1)},
\qquad \kappa=\frac{2\mu}{2\mu+1},\quad\alpha=\frac{r}{2r+d_Z}
\]
% This requirement is strictly weaker than the conclusion eq:phirate precisely
% when mu > 5/6, which is (P-f) and is the same inequality: rem:phirhyp(6).
\item[(P-a.3)] \emph{(source condition, stated \textbf{spectrally} and \textbf{at the compression})}
There are $\mu>0$ and a finite $\mathsf R_\ast$, neither depending on $n$, such that with
\begin{equation}\label{eq:compress}
\hat{\mathcal P}^{\,\mathcal S}_{\varphi}\;:=\;P_{\mathcal S_n}\,\hat{\mathcal A}^{\ast}_{\varphi}\hat{\mathcal A}_{\varphi}\big\vert_{\mathcal S_n}
\;:\;\mathcal S_n\to\mathcal S_n
\end{equation}
the \emph{compression} of $\hat{\mathcal A}^{\ast}_{\varphi}\hat{\mathcal A}_{\varphi}$ to
$\mathcal S_n$,
\[
P\Big(\ \bar v^{(k)}_\dagger=\big(\hat{\mathcal P}^{\,\mathcal S}_{\hat\varphi^{(k-1)}}\big)^{\mu}\varsigma^{(k)}
\ \text{ with }\ \big\|\varsigma^{(k)}\big\|_{L^{2}(\pi)}\le\mathsf R_\ast
\ \text{ for every }k\le\bar k\ \Big)\;\longrightarrow\;1 ,
\]
where $\bar v^{(k)}_\dagger:=P_{\mathcal S_n}\big(\vdagger-\hat\varphi^{(k-1)}\big)$ and $\bar k$ is the
stage count of \cref{def:lfstep}, \emph{which the proof's own construction puts at $\bar k\ge2$, so this
clause is \textbf{never} a one-stage condition and is always read at a $\bar v^{(k)}_\dagger$ formed at a
realised iterate as well as at one formed at the preliminary}.
\end{enumerate}
% THE bar-k >= 2 CLAUSE IS PRINTED HERE AS OF 2026-08-27, and the ground is the
% one that put the alpha restriction into this statement on 2026-08-05: a
% HYPOTHESIS a reader must act on should not require reading the proof.  bar-k is
% produced by the CONCLUSION and constructed in Step 4(a.2)-(b), where
% bar-k = bar-k_a + 1 with bar-k_a >= 1 by the "vee 1"; a reader checking (P-a.3)
% at the printed statement alone could not tell whether k = 1 was the only
% instance.  THE ANSWER DECIDES A GATE: at bar-k = 1 the clause is free of every
% iterate and a divergent stage-1 iteration would falsify eq:phirate; at
% bar-k >= 2 it is EMPTIED instead, and the proposition is vacuous rather than
% false on such a configuration.  Full disposition: rem:phirscope(7e).
% (P-a.3) IS NOT THE HILBERT-SCALE PAIR AND DOES NOT FOLLOW FROM IT, and the
% operator is named as the compression on purpose: rem:phirhyp(4), with the full
% disposition at rem:pasrc.
% ARCHAEOLOGY.  Rounds 1--11 assumed the Hilbert-scale pair; the replacement is a
% repair and not a strengthening, and is the most consequential amendment of that
% round.  Through round 11 the hypothesis named the restriction of A*A to S_n and
% the proof used the compression.
\item[(P-b)] \emph{(no hypothesis).} There is \textbf{no} hypothesis (P-b): the list runs
\textup{(P-a)}, \textup{(P-c)}--\textup{(P-g)};
% What the letter reserved -- an operator-estimation budget on a route Step 0 does
% not take -- is rem:phirhyp(5).
\item[(P-c)] \emph{(sieve approximation, a \textbf{lower} end for $L_x$)}
$\mathsf R\,L_x^{-s}=O\big(h_u\,\hat\delta_n\big)$;
% The factor h_u is not decoration, and (P-c) REPLACES A8's lower end rather than
% restating it: rem:phirhyp(7).
\item[(P-d)] \emph{(uniformity, an \emph{upper} end for $L_x$)} \suppeq{entropysieve}, so that
\suppref{Lemma}{unifop} clause (b) is available. \emph{(P-c) and (P-d) are the two ends of one window
and must be read together}: with $\hat\delta_n=n^{-\alpha}$ and $h_u\asymp\hat\delta_n^{1/8}$ they are
$L_x\gg n^{9\alpha/(8s)}$ and $L_x=o\big(n^{1-2\alpha}/\ln n\big)$, so the window is non-empty if and
only if
\begin{equation}\label{eq:lxwindow}
s\;>\;\frac{9\alpha}{8\,(1-2\alpha)} ;
\end{equation}
% What eq:lxwindow is numerically, and what carrying A8's lower end as well would
% have cost: rem:phirhyp(8).
\item[(P-e)] \emph{(stopping)} A7 in the amended form \eqref{eq:a7amend}, with
$\tau_{\mathrm{DP}}>2$;
\item[(P-f)] \emph{(the linearisation remainder is absorbed)}
$\delta_n^{2}=O\big(h_u^{2}\,\hat\delta_n\ln n\big)$, a hypothesis in which a bandwidth appears
explicitly, as (P-c) and (P-g) also do, and no claim of exhaustiveness is made in naming it so;
% What (P-f) reads at the unamended and at the amended bandwidth, and why the
% threshold is strict: rem:h6bandwidth, and rem:phirhyp(6).
\item[(P-g)] \emph{($x$-smoothing bias)} $\mathsf R\,h_x^{s-1/2}=O\big(h_u\,\hat\delta_n\big)$, an
\emph{upper} bound on $h_x$, which C5$'$ does not carry;
% THE FACTOR h_u IS NOT DECORATION HERE EITHER, AND IT WAS MISSING THROUGH
% DISPATCH 25.  lem:unifop Step 4 bounds the displacement of the two operator
% families in OPERATOR NORM by c_2 c_Pi h_u^{-1} max_i |Delta_{h_x} varphi_i -
% Delta varphi_i|; the h_u^{-1} is lem:intker(ii)'s ||partial_v C_{h_u}||_inf <=
% c_2 h_u^{-1}, the same h_u^{-1} that forces the net radius at lem:unifop(3c)
% and the h_u inside eq:entropy.  With eq:xsmoothbias that composes to
% c h_u^{-1} R h_x^{s-1/2}, so the step needs R h_x^{s-1/2} = O(h_u delta-hat_n).
% UNDER THE OLD FORM O(delta-hat_n) THE STEP DOES NOT CLOSE: at h_x at the top of
% its own admissible window the displacement is n^{alpha/8} delta-hat_n >>
% delta-hat_n, so (N1) does not follow, and (N1) feeds delta^eff, the rate
% recursion, eq:phirate and thm:rootn through (T-5).  (P-c) already carried the
% factor and prints "The factor h_u is not decoration"; (P-g) was the one place
% it was dropped, and rem:tccwhy(3) printed the correct arithmetic against it.
% WHAT MOVED WITH IT, all in one pass: lem:unifop's and lem:aerr's provisos,
% rem:tccwhy(3)'s quotation, rem:app5read(2) and eq:hxwindow, rem:phirhyp(1),
% rem:rootngroups(3) and (5), lem:scaleinv's cap inflation, and the budget-table
% row.  NO CONCLUSION IS LOST: the non-emptiness threshold rises from
% s > 1/2 + 8 alpha/(8 - alpha) to s > 1/2 + 9 alpha/(8 - alpha), whose supremum
% over alpha < 1/2 moves from 31/30 = 1.0333 to 11/10 = 1.1, and both running
% configurations of rem:h6bandwidth(5) clear the stricter threshold at s = 3
% with margins +2.06 and +2.13.  rem:app5read(2) prints that.
\end{enumerate}
and suppose finally
\begin{equation}\label{eq:muhalf}
\mu\;>\;\tfrac56 .
\end{equation}
Then there is a finite $\bar k\ge2$, not depending on $n$ and computable from $(\mu,\varsigma_0,\alpha)$,
such that the \textbf{infeasible} reading $\hat\varphi^{\pi}$ of \cref{def:lfstep} --- clause
\textup{(b)} of that definition with the adjoint formed at the population $\pi$ and the projection
$L^{2}(\pi)$-orthogonal onto $\mathcal S_n$, which is the object every step of the proof below runs, and
\textbf{not} the estimator $\hat\varphi$ that clause instructs, whose adjoint is formed at the floored
$\hat\pi$ --- satisfies
\begin{equation}\label{eq:phirate}
\big\|\hat\varphi^{\pi}-\vdagger\big\|_{L^{2}(\pi)}\;=\;O_P(\delta_n),
\qquad
\delta_n\;:=\;\big(\hat\delta_n\ln n\big)^{2\mu/(2\mu+1)} ,
\end{equation}
and $\delta_n=o_P(n^{-1/4})$ whenever \eqref{eq:frontier} holds.
\emph{What is \textbf{not} claimed is that $\hat\varphi$ satisfies \eqref{eq:phirate}}, and no hypothesis
of this proposition reaches that passage: it is the open item \suppref{Remark}{phirscope}\textup{(7)},
whose parts \textup{(7c)} and \textup{(7d)} are open. \emph{In particular this proposition does
\textbf{not} discharge A5(ii)}, which is stated at $\hat\varphi$; \cref{thm:rootn} carries that clause on
its own line, at \textup{(T-4)}, and carries it as a hypothesis.
% ############################################################################
% ##  SCOPED 2026-08-27.  THIS SENTENCE READ "the estimator of def:lfstep
% ##  satisfies eq:phirate", UNHEDGED, AND WAS A PRINTED SELF-CONTRADICTION.
% ##  def:lfstep(b) instructs the adjoint at the floored hat-pi; the proof below
% ##  says in terms that it runs the pair (S_n subset L^2(pi), L^2(Pi)), "which
% ##  is the infeasible reading", and that the feasible construction "is used at
% ##  no step of this proof"; and rem:phirscope(7) and rem:spacesread(4d) both
% ##  print, COMPILED, that the substitution is not established.  The statement
% ##  asserted what two compiled remarks deny.
% ##   THE ROUTE TAKEN IS SCOPING, DELIBERATELY, AND THE OTHER ROUTE WAS
% ##  EXAMINED AND REJECTED.  Adding the hat-pi hypothesis the corpus already
% ##  names -- ||hat-pi - pi||_{L^2(X,dx)} = O_P(delta_n), rem:phirscope(7c) --
% ##  closes (7c) and ONLY (7c).  (7d) is untouched by any rate on hat-pi: the
% ##  feasible recursion is governed by (pi/hat-pi) hat-A* hat-A, equivalent
% ##  renorming does not preserve range conditions, and (7d) records that the
% ##  perturbation route additionally needs a stability result for source
% ##  conditions under multiplicative perturbation which this paper does not
% ##  carry -- in a regime (eq:muhalf forces 2 mu > 5/3) where lem:srctransfer
% ##  (ii.a) REFUTES the interpolation transfer.  Keeping the unhedged claim with
% ##  that hypothesis added would contradict (7d) instead of (7c): the defect
% ##  would move, not close.
% ##   WHAT SURVIVES UNCHANGED: Steps 0-5 of the proof, every hypothesis, the
% ##  rate delta_n, eq:muhalf, eq:lxwindow and the budget table.  WHAT DOES NOT:
% ##  the claim to have verified A5(ii) primitively at the object def:crossfit(b)
% ##  plugs in.  thm:rootn is a CONDITIONAL theorem on that clause and (T-4)
% ##  already lists it.
% ############################################################################
% WHAT THE CONCLUSION DISCHARGES AND WHAT IT CONSUMES in the maintained block --
% A5(i) discharged, A5(ii) NOT (it is stated at hat-varphi and this conclusion is
% at hat-varphi^pi), A5(iii) never was, A5(vi') consumed, A5(vi) implied, A5(iv)
% and A5(v) untouched -- is rem:phirhyp(9).  [A5(ii)'s line moved 2026-08-27 with
% the scoping above; the comment said "A5(i)--(iii) discharged", which was already
% wrong about (iii) before the scoping and is now wrong about (ii) as well.]
% WHERE THE THRESHOLD eq:muhalf COMES FROM, and why 1/2 and not 5/6 is the
% coincidence worth naming: rem:phirhyp(10).
% WHAT HAPPENS AT 5/6, and why it is not a second coincidence: rem:phirhyp(11).
% WHY (P-a.3) IS THE LOAD-BEARING HYPOTHESIS and must be read in the same breath
% as eq:muhalf: rem:phirhyp(12).
% ARCHAEOLOGY.  The symbol delta_n^{(k)} that a round-9 draft used for the
% unobservable error is withdrawn; glossary C58.
% ARCHAEOLOGY.  A round-9 draft called the agreement at 5/6 a coincidence and said
% the two conditions were not restatements of one another; that sentence is struck.
\end{proposition}

\noindent\emph{Notes on \cref{prp:phirate}.} What its hypotheses carry, where the threshold
\eqref{eq:muhalf} comes from and why (P-a.3) is the load-bearing one are \suppref{Remark}{phirhyp}; why
(P-a.3) is spectral and at the compression is \suppref{Remark}{pasrc}; why (P-a)'s clauses are
w.p.$\to1$ envelopes is \suppref{Remark}{paenvelope}; the nonlinear route it declines is
\suppref{Remark}{lfnonlinear}; and what it does \emph{not} claim is \suppref{Remark}{phirscope}.
\emph{The last of those is not a pointer to fine print}: the conclusion above is scoped to the infeasible
reading $\hat\varphi^{\pi}$, and \suppref{Remark}{phirscope}\textup{(7)} is where the passage to the
estimator \cref{def:crossfit}\textup{(b)} actually plugs in is divided into five parts --- two settled
outright, one settled with an undischarged rider, two open.

% ===========================================================================
%  I.4  THE THEOREM
% ===========================================================================

\begin{theorem}[Asymptotic normality of the cross-fitted one-step estimator]\label{thm:rootn}
The one-step estimator is $\sqrt n$-consistent and asymptotically normal with the sandwich variance, under
mild ill-posedness and a propensity smooth enough to put the structural rate strictly inside $n^{-1/4}$.
% THE STATEMENT CARRIES CONDITIONS AND CONCLUSIONS AND NOTHING ELSE.  The
% hypothesis groups are tagged (T-1)--(T-5) and the additional hypotheses
% (R-a)--(R-e), on the model of prp:phirate's (P-a)--(P-g).  What each group
% does and does not carry is rem:rootngroups; how the conclusion is read is
% rem:rootnread.  Both live in results_supp.tex.
% ARCHAEOLOGY.  A draft of this theorem tagged the groups (H1)--(H5).  H1 is the
% Hilbert-scale link, it is cited INSIDE group (T-5), and a group label reading
% "H1" next to an assumption tag reading "H1" is the collision this file has
% repaired four times elsewhere.
Suppose
\begin{enumerate}
\item[(T-1)] \emph{(sampling, model, identification)} A0; A1; A1b(i), A1b(iii), A1b(iv);
\item[(T-2)] \emph{(the representer and the weight)} A2; A$2'$(i.a), A$2'$(i.b), A$2'$(iii); N1 at
$L_\star$, \emph{which is $p(\mathsf L-1)$ scalar conditions and \textbf{not} a condition on a
contraction of the $\mathsf L-1$ columns} (\cref{cor:centarity}), \emph{and which, together with
\textup{(R-a)}, pins the anchor A$2'$\textup{(iii)} leaves free}: \textup{(R-a)} forces
$\hat b^{(-k)}\to b_\star$, so $b_\star$ must be the N1 anchor of \suppref{Remark}{anchor} and the
construction must target it;
% ARCHAEOLOGY.  A3 stood in this group through dispatch 12.  No local rank
% condition on ell_star is in this group either; rem:rootngroups(1) says why.
\item[(T-3)] \emph{(index set, densities, propensity)} B1(i), (ii), (iv); C1(ii); C2$'$(i); A4(i),
A4(ii), A4(iii);
\item[(T-4)] \emph{(rates)} A5(i) and A5(ii) in the form \suppeq{deltalog}; A5(iii); and A5(vi) with
$s>\tfrac12$, read at each fold as in \cref{def:crossfit}(d). \emph{A5(ii) is \textbf{assumed} here and
is not delivered by \textup{(T-5)}}: it is stated at $\hat\varphi$, the estimator
\cref{def:crossfit}\textup{(b)} plugs in, while \cref{prp:phirate} concludes at the infeasible
$\hat\varphi^{\pi}$, and the passage between them is \suppref{Remark}{phirscope}\textup{(7)}, two of
whose five parts are open. \emph{This theorem is therefore conditional on that clause, and says so on its
own hypothesis line rather than in a remark};
% The group is NOT homogeneous in status, and it is LESS homogeneous from
% 2026-08-27 than it was: (T-5) delivers ONE of these clauses (A5(i), through
% lem:opnoise), no result in this paper delivers A5(iii), and A5(ii) -- delivered
% by (T-5) through dispatch 27 -- is now ASSUMED, prp:phirate having been scoped to
% hat-varphi^pi.  rem:rootngroups(2) records which is which.
\item[(T-5)] \emph{(the $\varphi$-step)} the hypotheses of \cref{prp:phirate}, read at the anchor
$\theta_\circ$ of \cref{def:crossfit}(b), H1; A5(vi$'$); A8's basis and upper-end clauses;
$p_K\ge s-\tfrac12$, which against \eqref{eq:lxwindow}'s floor gives $\alpha<\tfrac{68}{145}$; (P-a.1)--(P-a.3), (P-c)--(P-g); and
\eqref{eq:muhalf}, $\mu>\tfrac56$ strictly,
so that \eqref{eq:phirate} holds on each $\mathcal I_{-k}$ --- \emph{for the \textbf{infeasible} reading
$\hat\varphi^{\pi}$ of \cref{def:lfstep}, which is the only reading \cref{prp:phirate} delivers, and
\textbf{not} for the $\hat\varphi^{(-k)}$ that \cref{def:crossfit}\textup{(b)} plugs in}; so this group
supplies A5(i) and the fold reading of A5(vi$'$) and does \textbf{not} supply A5(ii), which
\textup{(T-4)} carries as a hypothesis;
\end{enumerate}
and suppose in addition
\begin{enumerate}
\item[(R-a)] \emph{(the weight is built as an anchored antiderivative, with a bounded derivative)} for each
$k$, $\hat L^{(-k)}=L^{\,\hat b^{(-k)}}_{\hat\ell^{(-k)}}$ in the sense of \cref{def:weight}(b), with
$\hat\ell^{(-k)}$ continuously differentiable in $u$ at every point and bounded,
\[
\max_{k\le\mathsf F}\Big\{\big\|\hat\ell^{(-k)}\big\|_\infty+\big\|\partial_u\hat\ell^{(-k)}\big\|_\infty\Big\}=O_P(1),
\qquad
\max_{k\le\mathsf F}E\big\|\hat b^{(-k)}(\Delta Z)\big\|^{4}=O_P(1),
\]
and $\max_{k\le\mathsf F}\|\hat L^{(-k)}-L_\star\|_{L^{2}(P)}=o_P(1)$;
% WHAT THIS CLAUSE DOES NOT REQUIRE -- sure boundedness of partial_u hat-ell --
% and what it adds to def:crossfit(c'.2)'s [0,+infty] typing: rem:rootngroups(3).
% ARCHAEOLOGY.  Dispatch 19 added a sure-boundedness conjunct on the ground that
% the two proof steps read finiteness straight off this hypothesis, and the same
% paragraph then recorded that "nothing was broken by the omission".  Both cannot
% be true of a minimal clause; the conjunct went and the O_P(1) display stayed, and
% the two proof sentences were re-grounded on the display rather than deleted.  A
% third site, thm:rootn Step 1, reads the clause only for the C^2 property of
% hat-L^{(-k)} and for the O_P rate, and is untouched either way.
\item[(R-b)] \emph{(fourth moments, and exactly two of them)} $E\|\Delta Z\|^{4}<\infty$ and
$E[\Delta U^{4}]<\infty$;
% Why a fourth moment of b_star is NOT needed: rem:rootngroups(4).
% ARCHAEOLOGY.  A round-10 draft carried a third clause, E||b_star||^4 < infinity.
\item[(R-c)] \emph{(the $\varrho\times L$ cross term)} \textbf{either} \textup{(c.1)} $D_k=\hat L^{(-k)}-L_\star$
satisfies N1$'$ for every $k$, \textbf{or} \textup{(c.2)}
$\max_k\|D_k\|_{L^{2}(P)}=O_P(\epsilon^{L}_n)$ with $r_n\,\epsilon^{L}_n=o(n^{-1/2})$;
\item[(R-d)] \emph{(no hypothesis).} There is \textbf{no} hypothesis (R-d): the list runs (R-a), (R-b),
(R-c), (R-e), on the \textup{(P-b)} pattern above;
\item[(R-e)] \emph{(the first-order bias)} \textbf{either} \textup{(e.1)} A5(iv) read at each fold, i.e.\ the
third display of \eqref{eq:foldrates} together with $\delta_n\epsilon_n=o_P(n^{-1/2})$, \textbf{or}
\textup{(e.2)} the hypotheses of \suppref{Lemma}{balance}\textup{(iii)}, the sieve solve sets term
\textup{(II)} of \suppeq{balance3} to zero, A8 holds with a B-spline basis of degree exceeding
$s-\tfrac12$, A5(vi) holds and \suppeq{nuhat} holds, together with
$\max_k\|\hat\ell^{(-k)}-\ell_\star\|_{L^{2}(P)}=o_P(1)$, \textbf{and} with
\begin{equation}\label{eq:re2window}
s\;>\;\frac{1}{2\,(1-2\alpha)} .
\end{equation}
% Why eq:re2window is stated on this branch and on no other, with the witness at
% which the branch would otherwise be empty: rem:rootngroups(5).
\end{enumerate}
Then, with $\hat\theta$ of \eqref{eq:onestep} and $J,\Sigma$ of \eqref{eq:JSig},
\begin{equation}\label{eq:rootn}
\begin{split}
\sqrt n\,\big(\hat\theta-\theta_\dagger\big)
&\;=\;J^{-1}\,\frac1{\sqrt n}\sum_{i=1}^{n}\psi_\star(O_i)\;+\;o_P(1)\\[2pt]
&\;\xrightarrow{\;d\;}\;\text{the centred Gaussian law on }\IR^{p}\text{ with covariance matrix }
J^{-1}\Sigma J^{-1\prime} ,
\end{split}
\end{equation}
and $\hat V\xrightarrow{p}J^{-1}\Sigma J^{-1\prime}$, so that for \textbf{every} $\vartheta\in\IR^{p}$ with
$\|\vartheta\|=1$,
\[
\sqrt n\,\vartheta'\big(\hat\theta-\theta_\dagger\big)\big/\big(\vartheta'\hat V\vartheta\big)^{1/2}
\;\xrightarrow{\;d\;}\;\text{the standard Gaussian law on }\IR ,
\]
the ratio being read on $\{\vartheta'\hat V\vartheta>0\}$, an event of probability tending to one.
% HOW THE CONCLUSION IS READ -- why the quantifier is "every" and not "every
% vartheta with vartheta' J^{-1} Sigma J^{-1'} vartheta > 0", why that is not a
% uniform lower bound, why the denominator can still vanish at finite n, what
% binds the letter vartheta, and why both limit laws are written out rather than
% given a symbol -- is rem:rootnread.
% ARCHAEOLOGY.  Through dispatch 13 the quantifier read "every vartheta with
% vartheta' J^{-1} Sigma J^{-1'} vartheta > 0"; the change is a strengthening of
% the statement at no cost in hypotheses, and the entailment was found by probing
% this very display.
% ARCHAEOLOGY.  Glossary C25.
% ARCHAEOLOGY (dispatch-23 change to how this theorem is PRINTED and to nothing it
% says).  Through dispatch 22 the two displays read N(0, J^{-1} Sigma J^{-1'}) and
% N(0,1) -- a law with two arguments on the letter lem:pfreeze's clause reserves.
% Step 5 of the proof already said "the limit law is written out and carries no
% symbol" while this statement printed the symbol twice.  The census behind glossary
% C84 named its own pattern as "N(", which does not match "N\big(", so the two
% displays here and the third at prp:medianpart were invisible to it: the proof was
% repaired and the three STATEMENTS were not.  Glossary C90 prints the per-file
% counts.
\end{theorem}

\noindent\emph{Notes on \cref{thm:rootn}.} What each hypothesis group does and does not carry is
\suppref{Remark}{rootngroups}; how the conclusion is read is \suppref{Remark}{rootnread}; what the theorem
covers and the six things it does not are \suppref{Remark}{rootnscope}; the clause-by-clause usage audit is
\suppref{Remark}{rootnhyp}; the logarithm in \eqref{eq:phirate} is \suppref{Remark}{rootnlog}.

\begin{corollary}[T2$'$: the rate floor, without \eqref{eq:frontier}]\label{cor:ratefloor}
The same decomposition delivers a rate whether or not the frontier holds; \emph{on the sieve branch, and
only there, every exponent in that rate is a function of the primitives}.
Suppose (T-1)--(T-3), (T-5), (R-a), (R-b) hold, together with A5(i), A5(vi) with $s>\tfrac12$,
$\delta_n\to0$, $\epsilon_n\to0$, $r_n\to0$, and (R-c.2), (R-e.1) with the two products left
\emph{unrestricted}. Then
\begin{equation}\label{eq:ratefloor}
\big\|\hat\theta-\theta_\dagger\big\|
\;=\;O_P\Big(n^{-1/2}\;+\;\delta_n\epsilon_n\;+\;\delta_n^{2}\;+\;r_n\delta_n\;+\;r_n\epsilon^{L}_n\Big).
\end{equation}
Under (R-e.2) in place of (R-e.1) the term $\delta_n\epsilon_n$ is replaced by
$O_P(n^{-1/2}\delta_n)+O_P(\mathsf R\,L_x^{-s})$, and no rate for $\hat\ell$ enters
\eqref{eq:ratefloor} at all. Finally, \eqref{eq:rootn} is the special case of \eqref{eq:ratefloor} in
which every term after the first is $o(n^{-1/2})$.
% HOW MUCH OF eq:ratefloor IS IN PRIMITIVES, stated by branch because the answer
% differs by branch, and the sharper display available on the sieve branch:
% rem:ratefloorprim.
% ARCHAEOLOGY.  A round-10 draft of this corollary claimed the stronger of the two
% answers on the weaker branch.
% ARCHAEOLOGY.  A round-10 draft listed a symbol "mu-prime", defined nowhere in
% this project and with no glossary row; it was deleted rather than defined.
\end{corollary}

\noindent\emph{Notes on \cref{cor:ratefloor}.} How much of \eqref{eq:ratefloor} is in primitives, and
which branch buys which answer, is \suppref{Remark}{ratefloorprim}.
